% Chapter 02.6: মোবাইল কমিউনিকেশন ও মোবাইল প্রজন্ম
\begingroup
\chTwoReadableTopicStyle
\section{মোবাইল কমিউনিকেশন ও মোবাইল প্রজন্ম}

\subsection{বাস্তব সমস্যা ও প্রয়োজন}

একজন শিক্ষার্থী বাসায় মোবাইল ডেটা দিয়ে অনলাইন ক্লাস করছে। সে শিক্ষকের কথা শুনছে, ভিডিও দেখছে এবং চ্যাট বক্সে প্রশ্ন পাঠাচ্ছে। ফোনটি বেতার সংকেত দিয়ে উপযুক্ত মোবাইল টাওয়ারে যুক্ত থাকে। এরপর মোবাইল নেটওয়ার্ক, সুইচিং কেন্দ্র ও ইন্টারনেট হয়ে অনলাইন ক্লাসের সার্ভারের সঙ্গে যোগাযোগ হয়।

এই যোগাযোগ চলার সময় শিক্ষার্থী স্থান বদলাতে পারে। চলমান অবস্থায়ও যোগাযোগ করতে পারা মোবাইল কমিউনিকেশনের প্রধান সুবিধা। একে \textbf{চলাচলের সুবিধা বা Mobility} বলে।

\begin{keypointbox}{মূল ধারণা}
মোবাইল যোগাযোগের মূল বিষয় ৫টি:
\begin{itemize}
\item মোবাইল ফোন বেতার সংকেত ব্যবহার করে।
\item মোবাইল টাওয়ার বা বেস স্টেশন নেটওয়ার্কে যুক্ত হওয়ার সুযোগ দেয়।
\item পুরো সেবার এলাকা ছোট ছোট সেলে ভাগ করা থাকে।
\item চলমান অবস্থায় ব্যবহারকারী এক সেল থেকে অন্য সেলে যেতে পারেন।
\item নেটওয়ার্কের প্রজন্ম অনুযায়ী গতি, সেবা ও বিলম্ব পরিবর্তিত হয়।
\end{itemize}
\end{keypointbox}

\subsection{মোবাইল কমিউনিকেশন}

\begin{definitionbox}{মোবাইল কমিউনিকেশন}
যে যোগাযোগব্যবস্থায় মোবাইল যন্ত্র বেতার সংকেত দিয়ে চলমান অবস্থায় কথোপকথন, ডেটা, লেখা, ছবি, অডিও, ভিডিও বা ইন্টারনেট সেবা আদান-প্রদান করতে পারে, তাকে মোবাইল কমিউনিকেশন বলে।
\end{definitionbox}

আধুনিক ICT-তে মোবাইল কমিউনিকেশন গুরুত্বপূর্ণ। ফোনে কথা বলা, SMS, MMS, মোবাইল ইন্টারনেট, ভিডিও কল, অনলাইন ব্যাংকিং, রাইড শেয়ারিং, GPS দিয়ে পথনির্দেশ, সামাজিক যোগাযোগ, ক্লাউড এবং IoT-ভিত্তিক সেবা এই ব্যবস্থায় ব্যবহার করা যায়।


\chTwoPair{\begin{keypointbox}{Mobile communication-এর বৈশিষ্ট্য}
\begin{itemize}
\item বেতার সংকেত ব্যবহার করে।
\item ব্যবহারকারী চলমান অবস্থায় যোগাযোগ করতে পারেন।
\item কথোপকথন ও ডেটা সেবা দেয়।
\item সেলুলার নেটওয়ার্কের গঠন ব্যবহার করে।
\item নেটওয়ার্কের সেবার আওতা অনুযায়ী সেবার মান বদলায়।
\item ব্যান্ডউইথ, সংকেতের শক্তি, বিলম্ব ও নেটওয়ার্কের ভিড় সেবার মানকে প্রভাবিত করে।
\end{itemize}
\end{keypointbox}}{\chTwoImage{chapters/Chapter02/images/mobile_communication_basic_model.png}{মোবাইল কমিউনিকেশনের সাধারণ মডেল}}



\subsection{Cellular Network-এর ধারণা}

\begin{definitionbox}{Cellular Network}
বড় ভৌগোলিক এলাকাকে ছোট ছোট সেলে ভাগ করে প্রতিটি সেলে বেস স্টেশন বা টাওয়ার দিয়ে মোবাইল সেবা দেওয়ার নেটওয়ার্ককে সেলুলার নেটওয়ার্ক বলে।
\end{definitionbox}

মোবাইল সেবা দেওয়ার জন্য বড় এলাকাকে ছোট ছোট সেলে ভাগ করা হয়। এই সরল মডেলে প্রতিটি সেলে একটি বেস স্টেশন বা BTS দেখানো হয়। ফোনটি উপযুক্ত সেবাদানকারী BTS-এর সঙ্গে রেডিও তরঙ্গ দিয়ে যোগাযোগ করে। পুরো এলাকা একটি টাওয়ার দিয়ে সরাসরি ঢেকে দেওয়াই এই ব্যবস্থার নকশা নয়।


\chTwoPair{\begin{keypointbox}{Cellular network-এর গুরুত্বপূর্ণ term}
\begin{itemize}
\item \textbf{Cell}: নির্দিষ্ট ছোট সেবার এলাকা।
\item \textbf{Base Station}: সেলের মোবাইল যন্ত্রগুলোর সঙ্গে বেতার যোগাযোগ করে।
\item \textbf{SIM}: গ্রাহকের পরিচয় যাচাইয়ে ব্যবহৃত হয়। SIM ও হ্যান্ডসেটের পরিচয়ের পার্থক্য পরের Multiple Access আলোচনায় আছে।
\end{itemize}
\end{keypointbox}}{\chTwoImage{chapters/Chapter02/images/cellular_network_cell_structure.png}{সেলুলার নেটওয়ার্কের সেল বিন্যাস}}


\subsection{Cellular Network কেন দরকার?}

বড় শহর বা দেশে মোবাইল সেবা দিতে সর্বত্র কার্যকর সংকেত প্রয়োজন। একটি শক্তিশালী টাওয়ার দিয়ে পুরো এলাকা ঢাকতে গেলে ফ্রিকোয়েন্সি পুনর্ব্যবহার, বহনক্ষমতা ও সেবার মানে সমস্যা হয়। ছোট সেলে ভাগ করলে দূরের সেলে একই ফ্রিকোয়েন্সি আবার ব্যবহার করা যায়। এতে বেশি ব্যবহারকারীকে সেবা দেওয়া সম্ভব হয়।

\begin{keypointbox}{Cellular network-এর সুবিধা}
\begin{itemize}
\item বড় এলাকায় সেবা দেওয়া যায়।
\item ফ্রিকোয়েন্সি পুনর্ব্যবহার করা যায়।
\item নেটওয়ার্কের বহনক্ষমতা বাড়ানো যায়।
\item বেশি ব্যবহারকারী একসঙ্গে সেবা পেতে পারেন।
\item চলাচলের সময় হ্যান্ডঅফ দিয়ে সংযোগ সচল রাখা যায়।
\item উপযুক্ত টাওয়ার পরিকল্পনায় সংকেতের মান বাড়ানো যায়।
\end{itemize}
\end{keypointbox}

\begin{mistakebox}{Mobile call path}
সাধারণ মোবাইল কলে একটি ফোন সরাসরি অন্য ফোনে কথা পাঠায় না। ফোন প্রথমে উপযুক্ত সেবাদানকারী টাওয়ার বা BTS-এ যুক্ত হয়। এরপর মোবাইল নেটওয়ার্ক কলের পথ নির্ধারণ করে।
\end{mistakebox}

\subsection{Mobile Network Architecture}

মোবাইল যোগাযোগে ফোন ও টাওয়ারের পাশাপাশি BTS, BSC, MSC, মূল নেটওয়ার্ক এবং ইন্টারনেট গেটওয়ের মতো অংশ কাজ করে। এই অংশগুলোর গঠন ও পারস্পরিক কাজ নিয়ে মোবাইল নেটওয়ার্কের আর্কিটেকচার তৈরি হয়।


\chTwoPair{\begin{keypointbox}{প্রধান অংশগুলোর কাজ}
\begin{itemize}
\item \textbf{Mobile Station}: ব্যবহারকারীর মোবাইল ফোন বা SIM-সমর্থিত যন্ত্র।
\item \textbf{BTS (Base Transceiver Station)}: মোবাইল ফোনের সঙ্গে বেতার সংকেত আদান-প্রদান করে।
\item \textbf{BSC (Base Station Controller)}: একাধিক BTS নিয়ন্ত্রণ করে এবং হ্যান্ডঅফের সমন্বয়ে সাহায্য করে।
\item \textbf{MSC (Mobile Switching Centre)}: কল স্থাপন, সুইচিং, পথ নির্ধারণ ও ব্যবহারকারীর চলাচল ব্যবস্থাপনা করে।
\item \textbf{Core Network}: ডেটা সেবা, পরিচয় যাচাই, বিল ও ইন্টারনেট সংযোগ পরিচালনা করে।
\end{itemize}
\end{keypointbox}}{\chTwoImage{chapters/Chapter02/images/mobile_network_architecture.png}{মোবাইল নেটওয়ার্ক আর্কিটেকচারের সরল রূপ}}


\subsection{Handoff বা Handover}

\begin{definitionbox}{Handoff বা Handover}
চলমান ব্যবহারকারী এক সেলের সেবার পরিসর থেকে অন্য সেলে গেলে কল বা ডেটা সংযোগ বজায় রেখে এক বেস স্টেশন থেকে অন্যটিতে সংযোগ স্থানান্তরের প্রক্রিয়াকে Handoff বা Handover বলে।
\end{definitionbox}

যেমন, একজন ব্যক্তি বাসে বসে ফোনে কথা বলছেন। বাস চলতে চলতে এক টাওয়ারের সেবার পরিসর থেকে অন্য টাওয়ারের পরিসরে যায়। কল সচল রাখতে নেটওয়ার্ক স্বয়ংক্রিয়ভাবে নতুন টাওয়ারে সংযোগ স্থানান্তর করে। এই স্থানান্তরই Handoff।


\chTwoPair{\begin{keypointbox}{Handoff-এর প্রয়োজনীয়তা}
\begin{itemize}
\item চলাচলের সময় কল বা ডেটা আদান-প্রদানের পর্ব সচল রাখতে সাহায্য করে।
\item ব্যবহারকারী এক সেল থেকে অন্য সেলে গেলে সংযোগ স্থানান্তর করে।
\item কল বিচ্ছিন্ন হওয়ার ঘটনা কমাতে সাহায্য করে।
\item মোবাইল যোগাযোগে চলাচলের সুবিধা বজায় রাখতে সাহায্য করে।
\item বাস, ট্রেন বা গাড়িতে দ্রুত চলাচলের সময় এটি গুরুত্বপূর্ণ।
\end{itemize}
\end{keypointbox}}{\chTwoImage{chapters/Chapter02/images/handoff_handover_process.png}{হ্যান্ডঅফ বা হ্যান্ডওভারের প্রক্রিয়া}}


\subsection{Roaming}

\begin{definitionbox}{Roaming}
নিজস্ব নেটওয়ার্কের সেবার পরিসরের বাইরে অন্য অপারেটর বা সহযোগী নেটওয়ার্ক ব্যবহার করে মোবাইল সেবা নেওয়ার সুবিধাকে Roaming বলে।
\end{definitionbox}

রোমিং দেশের মধ্যে বা আন্তর্জাতিক পর্যায়ে হতে পারে। অপারেটরের রোমিং চুক্তি থাকলে ব্যবহারকারী অন্য দেশে গিয়েও আন্তর্জাতিক রোমিং দিয়ে মোবাইল সেবা নিতে পারেন।


\chTwoPair{\begin{keypointbox}{Roaming-এর বৈশিষ্ট্য}
\begin{itemize}
\item নিজস্ব নেটওয়ার্কের বাইরে সেবা নেওয়া যায়।
\item দেশের মধ্যে বা আন্তর্জাতিক পর্যায়ে রোমিং হতে পারে।
\item অপারেটরগুলোর মধ্যে চুক্তি প্রয়োজন।
\item কল, SMS ও ডেটা সেবা ব্যবহার করা যায়।
\item রোমিংয়ের খরচ সাধারণত বেশি হতে পারে।
\end{itemize}
\end{keypointbox}}{\chTwoImage{chapters/Chapter02/images/mobile_roaming_concept.png}{রোমিংয়ের ধারণা}}


\subsection{মোবাইল প্রজন্ম বা Mobile Generation}

\begin{definitionbox}{Mobile Generation}
মোবাইল যোগাযোগ প্রযুক্তির ধারাবাহিক উন্নয়নের প্রতিটি ধাপকে মোবাইল প্রজন্ম বলে। যেমন—1G, 2G, 3G, 4G ও 5G।
\end{definitionbox}

\chTwoPair{প্রজন্ম শনাক্ত করতে সংকেতের রূপ, বেতার প্রযুক্তি, ডেটা সেবা, বহনক্ষমতা ও বিলম্ব বিবেচনা করতে হয়। একই অনলাইন কাজ একাধিক প্রজন্মে চলতে পারে। তাই একটি অ্যাপ্লিকেশনের নাম দিয়ে প্রজন্ম নিশ্চিত করা যায় না। নিচের ছকে প্রজন্মগুলোর পরিবর্তন দেখানো হয়েছে।}{\chTwoImage{chapters/Chapter02/images/mobile_generation_timeline.png}{মোবাইল প্রজন্মের ক্রমবিকাশ}}



\subsection{1G থেকে 5G: পরিবর্তন, প্রযুক্তি ও সেবা}
\begin{longtable}{|p{\dimexpr .12\linewidth-7pt\relax}|p{\dimexpr .31\linewidth-7pt\relax}|p{\dimexpr .57\linewidth-7pt\relax}|}\hline
\textbf{পর্যায়} & \textbf{Signal ও চেনার সূত্র} & \textbf{পরিবর্তন, ব্যবহার ও সীমা}\\\hline
\endfirsthead\hline\textbf{পর্যায়} & \textbf{Signal ও চেনার সূত্র} & \textbf{পরিবর্তন, ব্যবহার ও সীমা}\\\hline\endhead
1G & অ্যানালগ কথোপকথন; AMPS/NMT/TACS, FDMA & প্রাথমিক সেলুলার কথোপকথন; SMS ও ইন্টারনেট ছিল না। নয়েজ ও দুর্বল নিরাপত্তা; তুলনামূলক বড় যন্ত্র এবং ব্যাটারির কম কার্যদক্ষতা।\\\hline
2G & ডিজিটাল কথোপকথন; GSM/CDMA পরিবার & কথোপকথন ও SMS; উন্নত এনকোডিং, নিরাপত্তা ও বহনক্ষমতা। পরের উন্নয়নে MMS ও ডেটা যোগ হয়। সার্কিটভিত্তিক কলের পাশাপাশি প্যাকেট ডেটা আসে। বোতামযুক্ত ফোনে কল ও SMS পরিচিত ব্যবহার।\\\hline
2.5G / পরবর্তী 2G উন্নয়ন & GPRS; EDGE-কে প্রায়ই 2.75G বলা হয় & GPRS-এ প্যাকেট ডেটা এবং EDGE-এ উন্নত ডেটা হার পাওয়া যায়। মূল 2G থেকে 3G-তে যাওয়ার মধ্যবর্তী ধাপ; পৃথক পূর্ণ IMT প্রজন্ম নয়।\\\hline
3G & UMTS/WCDMA, HSPA; IMT-2000 & উন্নত মোবাইল ইন্টারনেট, ভিডিও কল, ই-মেইল ও মাল্টিমিডিয়া। স্মার্টফোনে ব্যাংকিং ও ই-কমার্স সহজ হয়। 2G-এর তুলনায় ডেটা হার বাড়ে।\\\hline
4G & LTE/LTE-Advanced; IP ব্রডব্যান্ড; IMT-Advanced & দ্রুত ইন্টারনেট, HD ভিডিও স্ট্রিমিং ও সম্মেলন, অনলাইন ক্লাস ও গেম। IP নেটওয়ার্কে VoLTE দিয়ে কল; 3G-এর তুলনায় কম বিলম্ব ও বেশি বহনক্ষমতা।\\\hline
5G & NR; IMT-2020 & বেশি বহনক্ষমতা, কম বিলম্ব ও বিপুল যন্ত্রের সংযোগ। স্মার্ট শহর, সেন্সর, স্বয়ংক্রিয় ব্যবস্থা, AR/VR ও উন্নত ব্রডব্যান্ডের চাহিদা পূরণের লক্ষ্য। অবকাঠামো ও সেবার আওতার সীমা রয়েছে।\\\hline
\end{longtable}
\textbf{ঐতিহাসিক নাম:} AMPS = Advanced Mobile Phone System; NMT = Nordic Mobile Telephone; TACS = Total Access Communication System। এগুলো 1G-এর পরিচিত ব্যবস্থা। জাপানে NTT-এর বাণিজ্যিক সেলুলার সেবা চালু হয় ১৯৭৯ সালে। যুক্তরাষ্ট্রে AMPS চালু হয় ১৯৮৩ সালে। দুটি ঘটনা পৃথক। D-AMPS (Digital AMPS) পরের ডিজিটাল ধারার নাম। নামের AMPS অংশ দেখে একে অ্যানালগ 1G বলা যায় না।

\subsection{5G-এর প্রয়োজনভিত্তিক ব্যবহার}
3GPP Release 15 থেকে প্রথম 5G মান নির্ধারণের ভিত্তি তৈরি হয়। NR বেতার সংযোগ এবং 5G Core সেবার প্রয়োজন অনুযায়ী নেটওয়ার্ক পরিচালনায় সাহায্য করে।
\chTwoPair{\begin{keypointbox}{বেশি Rate, কম Delay ও বেশি Device}
\textbf{eMBB:} উন্নত ব্রডব্যান্ড, 4K/8K ভিডিও, ক্লাউড গেমিং ও AR/VR।\par
\textbf{URLLC/Critical Communication:} কম বিলম্ব ও বেশি নির্ভরযোগ্যতা প্রয়োজন এমন শিল্প নিয়ন্ত্রণ, দূরনিয়ন্ত্রণ ও যানবাহনের সমন্বয়।\par
\textbf{Massive IoT:} কৃষির সেন্সর, স্মার্ট শহর ও পণ্য পরিবহনের গতিবিধি পর্যবেক্ষণে বহু ডিভাইসের সংযোগ।
\end{keypointbox}}{\chTwoImage{chapters/Chapter02/images/five_g_use_case_model..png}{প্রয়োজন অনুযায়ী 5G-এর ব্যবহার}}
\textbf{সহায়ক নকশা:} Network Slicing একই ভৌত নেটওয়ার্কে ভিন্ন সেবার জন্য পৃথক ভার্চুয়াল নেটওয়ার্কের ব্যবস্থা করে। Edge Computing-এ ব্যবহারকারীর কাছাকাছি ডেটা প্রক্রিয়াকরণ করা হয়। উন্নত অ্যান্টেনা, Beamforming ও বেশি সেলও ব্যবহার করা যায়। স্বয়ংচালিত গাড়ি বা দূর থেকে অস্ত্রোপচার চালাতে নেটওয়ার্কের পাশাপাশি নিরাপদ নিয়ন্ত্রণ, উপযুক্ত যন্ত্র ও পর্যাপ্ত নির্ভরযোগ্যতা প্রয়োজন।
\chTwoImage{chapters/Chapter02/images/mobile_generation_comparison_chart.png}{মোবাইল প্রজন্মের তুলনামূলক ধারণা}
\subsection{Mobile Communication-এর সীমাবদ্ধতা}

\begin{keypointbox}{সীমাবদ্ধতা}
\begin{itemize}
\item সংকেত দুর্বল হলে সেবার মান কমে।
\item নেটওয়ার্কে ভিড় বাড়লে গতি কমে যায়।
\item টাওয়ার থেকে দূরে গেলে সংকেত দুর্বল হতে পারে।
\item ভবন, আবহাওয়া ও বাধা সংকেতকে প্রভাবিত করতে পারে।
\item ব্যাটারির শক্তি বেশি ব্যয় হতে পারে।
\item ডেটার নিরাপত্তা ও ব্যক্তিগত গোপনীয়তার ঝুঁকি থাকে।
\item উচ্চগতির নেটওয়ার্ক স্থাপনের খরচ বেশি।
\end{itemize}
\end{keypointbox}

\subsection{Bandwidth, Latency ও Signal Strength-এর সম্পর্ক}

নতুন প্রজন্ম ব্যবহার করলেই উন্নত মোবাইল সেবা নিশ্চিত হয় না। ব্যান্ডউইথ, বিলম্ব, সংকেতের শক্তি, নেটওয়ার্কের ভিড়, টাওয়ারের দূরত্ব ও যন্ত্রের সক্ষমতার ওপর সেবার মান নির্ভর করে।

\begingroup
\begin{longtable}{|p{\dimexpr 0.22\linewidth-2\tabcolsep-\arrayrulewidth\relax}|p{\dimexpr 0.62\linewidth-2\tabcolsep-\arrayrulewidth\relax}|}
\hline
\textbf{বিষয়} & \textbf{প্রভাব} \\
\hline
\endfirsthead
\hline
\textbf{বিষয়} & \textbf{প্রভাব} \\
\hline
\endhead

Bandwidth & ব্যান্ডউইথ বেশি হলে একই সময়ে বেশি ডেটা পাঠানোর সুযোগ থাকে। \\
\hline
Latency & বিলম্ব কম হলে ভিডিও কল, অনলাইন গেম ও সরাসরি ক্লাসের মান উন্নত হতে পারে। \\
\hline
Signal Strength & শক্তিশালী সংকেত সংযোগ স্থিতিশীল রাখতে সাহায্য করে। \\
\hline
Congestion & ব্যবহারকারী বেশি হলে গতি কমতে পারে। \\
\hline
Tower Distance & টাওয়ার দূরে হলে সংকেত দুর্বল হতে পারে। \\
\hline
Device Capability & পুরোনো যন্ত্র উচ্চগতির নেটওয়ার্কের সুবিধা পুরোপুরি ব্যবহার করতে নাও পারে। \\
\hline
\end{longtable}
\endgroup

\begin{mistakebox}{4G/5G speed}
ফোনে 4G/5G লেখা থাকলেই ইন্টারনেট সবসময় দ্রুত হয় না। সংকেতের শক্তি, ব্যান্ডউইথ, বিলম্ব, নেটওয়ার্কের ভিড়, যন্ত্রের মান এবং সার্ভারের সাড়া দেওয়ার সময় গতিকে প্রভাবিত করে।
\end{mistakebox}

\subsection{Mobile Communication-এর ব্যবহার}

\begingroup
\begin{longtable}{|p{\dimexpr 0.25\linewidth-2\tabcolsep-\arrayrulewidth\relax}|p{\dimexpr 0.55\linewidth-2\tabcolsep-\arrayrulewidth\relax}|}
\hline
\textbf{ক্ষেত্র} & \textbf{ব্যবহার} \\
\hline
\endfirsthead
\hline
\textbf{ক্ষেত্র} & \textbf{ব্যবহার} \\
\hline
\endhead

শিক্ষা & অনলাইন ক্লাস, ডিজিটাল উপকরণ, ই-লার্নিং \\
\hline
স্বাস্থ্য & টেলিমেডিসিন, জরুরি কল, স্বাস্থ্য পর্যবেক্ষণ \\
\hline
ব্যবসা & মোবাইল ব্যাংকিং, ই-কমার্স, অনলাইন লেনদেন \\
\hline
পরিবহন & GPS দিয়ে পথনির্দেশ, রাইড শেয়ারিং, যানবাহনের অবস্থান জানা \\
\hline
বিনোদন & ভিডিও স্ট্রিমিং, অনলাইন গেম, সামাজিক যোগাযোগ \\
\hline
নিরাপত্তা & জরুরি সতর্কবার্তা, অবস্থান জানা, নজরদারিতে সহায়তা \\
\hline
স্মার্ট ব্যবস্থা & IoT, স্মার্ট বাসা, স্মার্ট শহর, স্বয়ংক্রিয় ব্যবস্থা \\
\hline
\end{longtable}
\endgroup

\subsection{Cell-এর আকার, Frequency Reuse ও Tower নির্বাচন}
পাহাড়, ভবন, অ্যান্টেনা ও শক্তির প্রভাবে বাস্তব সেলের সেবার পরিসর অনিয়মিত হয়। পরিকল্পনার চিত্রে ষড়ভুজ ব্যবহার করা হয়। ষড়ভুজ প্রায় বৃত্তাকার পরিসর বোঝায় এবং পাশাপাশি বসিয়ে ফাঁক ছাড়াই পুরো এলাকা দেখানো যায়। একটি বিশাল টাওয়ারের বদলে ছোট সেল ব্যবহার করলে ফ্রিকোয়েন্সি পুনর্ব্যবহার এবং বেশি ব্যবহারকারীকে সেবা দেওয়া সহজ হয়।

\textbf{ফ্রিকোয়েন্সি পুনর্ব্যবহার:} যথেষ্ট দূরের সেলে একই ফ্রিকোয়েন্সি আবার ব্যবহার করা যায়। কাছের সেলগুলোর বিঘ্ন নিয়ন্ত্রণ করতে পরিকল্পনা প্রয়োজন। মোবাইল সাধারণত সবচেয়ে উপযুক্ত সেবাদানকারী সেল বেছে নেয়। সেটি সবসময় সবচেয়ে কাছের টাওয়ার নয়। সংকেতের মান, কাজের চাপ ও নেটওয়ার্কের সিদ্ধান্তও নির্বাচনে ভূমিকা রাখে।

\textbf{GSM/CDMA-এর সেবার পরিসরের উদাহরণ:} প্রচলিত তুলনায় GSM-এর প্রায় ৩৫ km এবং CDMA-এর প্রায় ১১০ km সেলের পরিসর ধরা হয়। GSM-এর প্রচলিত Timing Advance-এর সীমা প্রথম সংখ্যাটির একটি প্রযুক্তিগত প্রেক্ষাপট। বিশেষ সম্প্রসারণে কার্যকর দূরত্ব ভিন্ন হতে পারে। CDMA-র সংখ্যাটি নির্দিষ্ট ঐতিহাসিক নকশার উদাহরণ; সব CDMA সেলের আবশ্যিক পরিসর নয়। ছোট শহুরে সেল বা ঘরের ভেতরের সেলও হতে পারে। শুধু প্রযুক্তির নাম দিয়ে বাস্তব সেলের ব্যাসার্ধ জানা যায় না।

\textbf{GSM-900-এর ফ্রিকোয়েন্সির হিসাব:} প্রচলিত GSM-900-এ ফোন থেকে বেস স্টেশনে Uplink হলো ৮৯০–৯১৫ MHz। বিপরীত পথে Downlink হলো ৯৩৫–৯৬০ MHz। প্রতিটি দিকের ব্যান্ডউইথ $915-890=25$ MHz। সমতুল্য Uplink ও Downlink ক্যারিয়ারের ব্যবধান ৪৫ MHz। রেডিও চ্যানেলের ব্যবধান ২০০ kHz এবং একটি TDMA ফ্রেমে ৮টি সময়ের ভাগ থাকে। GSM-এর অন্য ব্যান্ড ও সম্প্রসারিত ব্যান্ডে ফ্রিকোয়েন্সি আলাদা। এই সংখ্যা সব মোবাইল প্রজন্মের জন্য প্রযোজ্য নয়।

\subsection{একই মাধ্যম বহু গ্রাহক ভাগ করে কীভাবে?}
\textbf{GSM-এ Frequency Hopping:} নির্ধারিত ক্রম অনুসরণ করে একাধিক অনুমোদিত ক্যারিয়ার ফ্রিকোয়েন্সিতে পালাক্রমে সংকেত পাঠানোকে Frequency Hopping বলে। প্রেরক ও প্রাপক একই ক্রম অনুসরণ করে। শুধু বেশি ফ্রিকোয়েন্সি বেছে নেওয়াই এর কাজ নয়। নির্দিষ্ট ফ্রিকোয়েন্সির বিঘ্ন বা সংকেত ম্লান হওয়ার প্রভাব বিভিন্ন ফ্রিকোয়েন্সিতে ছড়িয়ে দিতে এই পদ্ধতি সাহায্য করতে পারে। Frequency Hopping এবং এক সেল থেকে অন্য সেলে Handoff পৃথক ঘটনা।

FDMA = Frequency Division Multiple Access; TDMA = Time Division Multiple Access; CDMA = Code Division Multiple Access। তিন পদ্ধতিতে ভাগ করার ভিত্তি যথাক্রমে ফ্রিকোয়েন্সি, সময় ও কোড।
\begin{tabularx}{\linewidth}{|p{.18\linewidth}|X|X|}
\hline
\textbf{পদ্ধতি} & \textbf{ভাগ করার ভিত্তি} & \textbf{সহজ উদাহরণ ও প্রজন্ম}\\\hline
FDMA & পৃথক ফ্রিকোয়েন্সি চ্যানেল & একই সময়ে পৃথক ফ্রিকোয়েন্সিতে যোগাযোগ; 1G-এর পরিচিত পদ্ধতি।\\\hline
TDMA & একই চ্যানেলে পৃথক সময়ের ভাগ & পালাক্রমে ছোট সময়ের অংশ বরাদ্দ; 2G GSM-এ ব্যবহৃত।\\\hline
CDMA & পৃথক কোড & কোড অনুযায়ী তথ্য আলাদা করা; 2G-এর কিছু প্রযুক্তি পরিবার ও 3G WCDMA।\\\hline
\end{tabularx}
একই মাধ্যম বহু ব্যবহারকারী ভাগ করে ব্যবহারের ধারণা হলো Multiple Access। TDMA-তে সময় ভাগ করা হলেও ব্যবহারকারীর পুরো টেলিফোন কথোপকথন হাফ-ডুপ্লেক্স হয়ে যায় না। সংযোগের দিক ও বেতার সম্পদ ভাগ করার পদ্ধতি আলাদা বিষয়।

\textbf{GSM ও CDMA-র পার্থক্য:} প্রচলিত GSM-এ প্রতিটি ফ্রিকোয়েন্সি চ্যানেলের মধ্যে পৃথক সময়ের ভাগ বরাদ্দ করা হয়। অর্থাৎ FDMA ও TDMA দুটিই কাজে লাগে। CDMA-তে পৃথক কোড দিয়ে একই যোগাযোগসম্পদে ব্যবহারকারীদের তথ্য আলাদা করা হয়। কোড ব্যবহারকারীর ব্যক্তিগত ফোন নম্বর নয়। এই ব্যবস্থায় উন্নত শক্তি নিয়ন্ত্রণ ও উপযুক্ত গ্রাহক যন্ত্র প্রয়োজন। সব CDMA ব্যবস্থার ব্যাটারির ব্যবহার, কার্যকর দূরত্ব বা কথার মান সমান নয়।
\chTwoPair{\textbf{Subscriber-এর পরিচয়:} SIM = Subscriber Identity Module। গ্রাহকের পরিচয় ও তা যাচাইয়ের তথ্য অপারেটরের সেবা নিতে সাহায্য করে। IMSI (International Mobile Subscriber Identity) গ্রাহকের স্থায়ী এবং TMSI (Temporary Mobile Subscriber Identity) অস্থায়ী পরিচয়। IMEI (International Mobile Equipment Identity) সাধারণত ১৫ অঙ্কের হ্যান্ডসেট পরিচয়। এটি গ্রাহকের পরিচয় বা ফোন নম্বর নয়।}{\textbf{Handset বদলানোর শর্ত:} প্রচলিত GSM-এ অন্য সমর্থিত হ্যান্ডসেটে SIM ব্যবহার করা যায়। বেতার ব্যান্ড, প্রযুক্তি, অপারেটরের লক ও অনুমতি সামঞ্জস্যপূর্ণ হতে হবে। কিছু CDMA হ্যান্ডসেট অপারেটর সক্রিয় করে দেয়। কিছু হ্যান্ডসেটে RUIM (Removable User Identity Module) থাকে। তাই সব CDMA ব্যবস্থায় কার্ড বা রোমিং নেই বলা ঠিক নয়।\par
\textbf{পুরোনো পরিচয়গুলোর নাম:} ESN (Electronic Serial Number) যন্ত্রের, MIN (Mobile Identification Number) গ্রাহকের এবং SID (System Identification) সিস্টেমের পরিচয় বোঝায়। ঐতিহাসিক সেলুলার বা CDMA ব্যবস্থা বুঝতে নামগুলো প্রয়োজন। সব নতুন নেটওয়ার্কে একই পরিচয়সূচক ব্যবহৃত হয় না।}

\subsection{প্রযুক্তির পূর্ণরূপ, চালুর সাল ও Rate}
\begin{multicols}{2}
\textbf{GSM, GPRS ও EDGE:} GSM: Global System for Mobile Communications। GPRS: General Packet Radio Service। EDGE: Enhanced Data rates for GSM Evolution। GPRS ও EDGE দিয়ে GSM-ভিত্তিক নেটওয়ার্কে উন্নত প্যাকেট ডেটা সেবা যুক্ত হয়।\par
\textbf{UMTS, LTE ও NR:} UMTS: Universal Mobile Telecommunications System। LTE: Long Term Evolution। NR: New Radio। পরীক্ষার উদ্দীপকে LTE দিয়ে সাধারণত 4G বোঝানো হয়। LTE-Advanced কঠোর IMT-Advanced লক্ষ্য পূরণের দিকে উন্নত রূপ।\par
\end{multicols}

\textbf{কয়েকটি নির্দিষ্ট চালুর সাল:} প্রথম 1G নেটওয়ার্কের ১৯৭৯/১৯৮৩ সালের প্রেক্ষাপট আগের 1G অংশে আছে। Finland-এ Nokia-র যন্ত্রে প্রথম আনুষ্ঠানিক GSM কল হয় ১ জুলাই ১৯৯১ সালে। TeliaSonera-র Stockholm LTE সেবায় ২০০৯ সালের ডিসেম্বরে Ericsson-এর নেটওয়ার্ক ব্যবহৃত হয়। সেবাটি বাজারে 4G নামে পরিচিত। একটি দেশে সেবা উদ্বোধনের সাল পুরো বিশ্বের প্রজন্ম বদলের সাল নয়।

\textbf{5G-এর চালুর সাল:} ২০১৯ সালেই প্রাথমিক বাণিজ্যিক 5G সেবা চালু ছিল। IMT-2020 নামের মধ্যে ২০২০ থাকলেও সব দেশে প্রথম 5G চালুর সাল ২০২০ নয়। প্রযুক্তির উন্নয়ন, মান অনুমোদন ও স্থানীয় সেবা চালুর তারিখ পৃথক হতে পারে।

\textbf{FOMA:} Freedom Of Mobile multimedia Access হলো NTT DOCOMO-র ২০০১ সালে চালু করা 3G সেবার ঐতিহাসিক নাম। এটি পুরো 3G প্রজন্মের পূর্ণরূপ নয়। WCDMA-তে W দিয়ে Wideband বোঝায়।

\textbf{মান ও প্রযুক্তির নাম:} IMT = International Mobile Telecommunications। ITU-র IMT-2000/Advanced/2020 যথাক্রমে 3G/4G/5G পরিচয়ের সঙ্গে সম্পর্কিত মানের পরিবার। 3GPP = Third Generation Partnership Project। এটি মান উন্নয়নের অংশীদারত্ব; ফোন বা একক বেতার লিংকের নাম নয়। HSPA = High Speed Packet Access। HSDPA (High Speed Downlink Packet Access) ডাউনলোড এবং HSUPA (High Speed Uplink Packet Access) আপলোডের উন্নয়নের পরিচিত নাম। RAT = Radio Access Technology; GSM বা LTE-এর মতো বেতার সংযোগের প্রযুক্তি বোঝায়।

\textbf{3G-এর প্রযুক্তি ও 4G-এর সাবক্যারিয়ার:} 3G-তে UMTS/WCDMA ছাড়াও CDMA2000 পরিবার আছে। TD-CDMA ও TD-SCDMA (Time Division Synchronous CDMA)-ও এই প্রজন্মের প্রযুক্তি। একটি নাম পুরো 3G-এর একমাত্র পদ্ধতি নয়। OFDM (Orthogonal Frequency Division Multiplexing)-এ ডেটা সমন্বিত বহু সাবক্যারিয়ারে ভাগ করে বহন করা হয়। OFDMA-তে সাবক্যারিয়ারের সম্পদ বিভিন্ন ব্যবহারকারীকে বরাদ্দ দেওয়া হয়। LTE-এর Downlink এবং Wi-Fi/WiMAX-এর উপযুক্ত সংস্করণে এসব ব্যবহৃত হয়। এটি পৃথক তামার তারে Parallel Transmission-এর একই ধারণা নয়।

\textbf{3G-এর প্রাথমিক হারের উদাহরণ:} পুরোনো IMT-2000 পাঠে দ্রুত চলমান অবস্থায় প্রায় ১৪৪ kbps, হাঁটা বা ধীর চলাচলে ৩৮৪ kbps এবং স্থির বা ঘরের ভেতরে ২ Mbps-এর লক্ষ্য পাওয়া যায়। পরে HSPA-তে ডেটা হার বাড়ে। তাই 3G-এর গতি সবসময় শুধু ২ Mbps নয়। এই মানগুলো ব্যবহারে চলাচল, বেতার প্রযুক্তি ও পরিবেশের শর্ত বিবেচনা করতে হয়।

\textbf{ডেটা হারের শর্ত:} 4G-তে ১০০ Mbps (দ্রুত চলাচল) ও ১ Gbps (ধীর চলাচল) পরিচিত গবেষণা লক্ষ্য। IMT-2020/5G-এর আদর্শ সর্বোচ্চ Downlink ২০ Gbps এবং Uplink ১০ Gbps। প্রতি বর্গকিলোমিটারে প্রায় ১০ লক্ষ যন্ত্র সংযোগের লক্ষ্যও আছে। এগুলো নির্দিষ্ট পরীক্ষার শর্তের লক্ষ্য। প্রত্যেক ফোনের নিশ্চিত গতি বা সব লক্ষ্য একই সঙ্গে পাওয়ার প্রতিশ্রুতি নয়। পরিষেবার প্রচলিত হারের ছক অংশ \ref{ch2:c01-12}-এ আছে।

\textbf{WWWW-এর পরিচয়:} World Wide Wireless Web দিয়ে 5G-এর বিস্তৃত বেতার সেবার ধারণা বোঝানো হয়। Wireless World Wide Web শব্দক্রমও প্রচলিত। এটি IMT-2020 বা NR-এর মতো আনুষ্ঠানিক বেতার মান নয়। IoT, কম বিলম্ব, বেশি বহনক্ষমতা ও 4K ভিডিওর বৈশিষ্ট্যগুলো একত্রে 5G-এর প্রসঙ্গ নির্দেশ করতে পারে। শুধু 4K ভিডিও অন্য উপযুক্ত ব্রডব্যান্ডেও চলে। তাই ভিডিওর রেজল্যুশন একা প্রজন্ম শনাক্ত করে না।

\textbf{আর্কিটেকচারের সীমা:} BTS–BSC–MSC-সহ চিত্রটি প্রচলিত GSM/2G-এর সরল মডেল। 4G/5G-এর বেতার অংশ ও মূল নেটওয়ার্কে নাম এবং দায়িত্ব ভিন্ন হয়। নতুন প্রযুক্তির সম্ভাবনা ব্যাখ্যা করা যায়। তবে শুধু 5G থাকলেই দূর থেকে অস্ত্রোপচার বা স্বয়ংচালিত গাড়ির সব ব্যবস্থা নিরাপদে চলবে—এমন সিদ্ধান্ত দেওয়া যায় না।

\begin{keypointbox}{MIMO: একাধিক Antenna কেন?}
MIMO (Multiple Input Multiple Output)-তে প্রেরক ও প্রাপক প্রান্তে একাধিক অ্যান্টেনা ব্যবহার করা হয়। উপযুক্ত পরিবেশে একই সময় ও ফ্রিকোয়েন্সির সম্পদে একাধিক ডেটা প্রবাহ বহন করা যায়। এতে সংযোগের বহনক্ষমতা বাড়তে পারে। LTE/4G ও আধুনিক Wi-Fi-তে এই ধারণা ব্যবহৃত হয়। রাস্তার একাধিক লেন দিয়ে পাশাপাশি গাড়ি যাওয়ার সঙ্গে এর তুলনা করা যায়। তবে শুধু অ্যান্টেনা বাড়ালেই সব পরিবেশে একই গতি পাওয়া যায় না। সংকেত, চ্যানেল ও যন্ত্রের সক্ষমতাও প্রয়োজন।
\end{keypointbox}


% Original topic-level problem, solution and MCQ reasoning.
\subsection{বাস্তব সমস্যা: চলন্ত Bus-এ Call ও Data}
\chTwoPair{\textbf{সমস্যা:} বাস চলার সময় ফোনে কথা চলছে। পরে অন্য অপারেটরের সহযোগী নেটওয়ার্কে সেবা নিতে হলো। একই ফোনে ভিডিও কখনো দ্রুত চলে, কখনো আটকে যায়।\par
\textbf{সমাধান:} সেবাদানকারী সেল বদলানো হলো Handoff। নিজস্ব নেটওয়ার্কের বাইরে সহযোগীর সেবা নেওয়া হলো Roaming। ভিডিও চলার মান সংকেত, বহনক্ষমতা, ভিড় ও বিলম্ব অনুযায়ী বদলায়। শুধু 4G লেখা থাকলেই নির্দিষ্ট ডেটা হার নিশ্চিত হয় না।}{\textbf{Mobile Link-এর Uplink ও Downlink:} ফোন থেকে স্টেশনে পাঠানো হলো Uplink; স্টেশন থেকে ফোনে পাঠানো হলো Downlink। স্যাটেলাইটে ভূকেন্দ্র থেকে উপগ্রহের নামের সঙ্গে কোন প্রান্তকে ভিত্তি ধরা হচ্ছে, তা মিলিয়ে বুঝতে হয়।\par
\textbf{Packet Data ও Generation:} GPRS/EDGE প্রচলিত 2G-এর ডেটা সেবার উন্নত রূপ। প্যাকেট ডেটা একটিমাত্র প্রজন্মে সীমাবদ্ধ নয়।\par}
\begin{multicols}{2}
\textbf{SMS, MMS ও Billing:} SMS-এর পূর্ণরূপ Short Message Service এবং MMS-এর পূর্ণরূপ Multimedia Messaging Service। SMS দিয়ে লেখা এবং MMS দিয়ে ছবি বা অডিওসহ বার্তা পাঠানো যায়। প্রচলিত প্রশ্নে এগুলো ডিজিটাল বা 2G ধারার পরিচয় হিসেবে থাকতে পারে। Prepaid হলো আগে অর্থ দিয়ে সেবা ব্যবহারের বিল পদ্ধতি। প্রচলিত প্রজন্মের ছকে এটি 2G-এর সঙ্গে দেখানো হয়।\par
\textbf{Circuit ও Packet Switching:} Circuit Switching-এ একটি যোগাযোগ পর্বের জন্য নির্দিষ্ট যোগাযোগসম্পদ সংরক্ষণ করা হয়। Packet Switching-এ ডেটা ছোট প্যাকেটে ভাগ করে নেটওয়ার্কের সম্পদ ভাগাভাগি করা হয়। GPRS-এও প্যাকেট ডেটা আছে। 3G-তে এর ব্যবহার বিস্তৃত হয় এবং 4G-তে IP-ভিত্তিক ব্যবস্থা প্রধান।\par
\textbf{Green Phone-এর পরিচয়:} পরীক্ষায় Green Phone নামটি CDMA-এর প্রচলিত পরিচয় হিসেবে থাকতে পারে। CDMA-র শক্তি নিয়ন্ত্রণ প্রয়োজনমতো প্রেরণশক্তি রেখে বিঘ্ন ও ব্যাটারির ব্যবহার কমাতে সাহায্য করে। Green Phone আধুনিক মান বা প্রত্যয়নপত্রের নাম নয়।\par
\textbf{Portable Cellular Phone-এর ইতিহাস:} Motorola-র Martin Cooper ৩ এপ্রিল ১৯৭৩ সালে হাতে বহনযোগ্য সেলুলার ফোনে প্রকাশ্যে কল করেন। DynaTAC 8000X বাণিজ্যিক বহনযোগ্য ফোন হিসেবে আসে ১৯৮৩ সালে। প্রথম প্রদর্শন ও বাজারে আসার সাল পৃথক। Cooper-এর অবদান সেলুলার হ্যান্ডসেটের ইতিহাসে গুরুত্বপূর্ণ; তিনি সব বেতার যোগাযোগের একমাত্র উদ্ভাবক নন।\par
\end{multicols}
\textbf{পরিষেবার Rate:} VoIP, Streaming Video ও HD Video Conference-এর প্রচলিত সংখ্যাগত উদাহরণ যোগাযোগ ও ব্যান্ডউইথের অংশে দেওয়া আছে (অংশ \ref{ch2:c01-12}); সেখানকার পরিষেবা-ছক অনুসরণ করবে।

\textbf{অবস্থান ও ডেটা সেবা:} GPS = Global Positioning System। এটি অবস্থান নির্ণয়ের স্যাটেলাইট ব্যবস্থা। GPRS-এর মোবাইল ডেটা সেবা থেকে GPS আলাদা। IP টেলিফোনি, গেমিং এবং HD/3D ভিডিও মোবাইল ব্রডব্যান্ডের ব্যবহার। এসব কাজে প্রয়োজনীয় ডেটা হার, যন্ত্র ও কোডেক সামঞ্জস্যপূর্ণ হতে হয়।


\subsection{সাধারণ ভুল ও সংশোধন}

\smallhead{১. 1G ও 2G শুধু speed দিয়ে আলাদা করা}
1G ও 2G-এর মূল পার্থক্য সংকেতের ধরনে। 1G-তে অ্যানালগ এবং 2G-তে ডিজিটাল যোগাযোগ ব্যবহৃত হয়। এই পার্থক্য বাদ দিয়ে শুধু গতি উল্লেখ করলে ব্যাখ্যা অসম্পূর্ণ থাকে।

\smallhead{২. Handoff ও roaming এক মনে করা}
Handoff-এ চলমান ব্যবহারকারীর সংযোগ এক টাওয়ার থেকে অন্য টাওয়ারে স্থানান্তরিত হয়। Roaming-এ নিজস্ব নেটওয়ার্কের বাইরে অন্য নেটওয়ার্কের সেবা নেওয়া হয়। দুটি প্রক্রিয়া আলাদা।

\smallhead{৩. 3G মানেই শুধু video call মনে করা}
3G শুধু ভিডিও কলের জন্য নয়। 3G-তে মোবাইল ইন্টারনেট, মাল্টিমিডিয়া, ই-মেইল, ব্রাউজিং ও ডেটা সেবাও উন্নত হয়।

\smallhead{৪. 5G শুধু বেশি speed মনে করা}
5G-তে বেশি গতি গুরুত্বপূর্ণ। কম বিলম্ব, বিপুল যন্ত্র সংযোগ এবং IoT সেবার সমর্থনও 5G-এর উল্লেখযোগ্য বৈশিষ্ট্য।

\subsection{MCQ অনুশীলন}

\par
\textbf{MCQ}
\begin{multicols}{2}\begin{enumerate}
\item\chTwoMCQBlock{মোবাইল যোগাযোগের প্রধান বৈশিষ্ট্য কোনটি?
\par ক. শুধু তারযুক্ত সংযোগ \quad খ. চলাচলের সময় যোগাযোগের সুযোগ \par গ. শুধু প্রিন্টার ভাগাভাগি \quad ঘ. শুধু ডেটা সংরক্ষণ }
\item\chTwoMCQBlock{সেলুলার নেটওয়ার্কে বড় এলাকাকে কী হিসেবে ভাগ করা হয়?
\par ক. ফাইল \quad খ. সেল \par গ. ফোল্ডার \quad ঘ. প্যাকেট }
\item\chTwoMCQBlock{BTS-এর কাজ কী?
\par ক. শুধু মুদ্রণ \par খ. মোবাইল ফোনের সঙ্গে রেডিও সংকেত আদান-প্রদান \par গ. ডেটাবেজের তথ্য সাজানো \quad ঘ. ছবি সম্পাদনা }
\item\chTwoMCQBlock{1G-তে যোগাযোগের প্রধান ধরন কী ছিল?
\par ক. শুধু ডিজিটাল SMS \quad খ. অ্যানালগ ভয়েস \par গ. 5G IoT \quad ঘ. ফাইবার অপটিক }
\item\chTwoMCQBlock{2G-এর গুরুত্বপূর্ণ সেবার উদাহরণ কোনটি?
\par ক. SMS \quad খ. শুধু ভার্চুয়াল রিয়েলিটি \par গ. কোয়ান্টাম কম্পিউটিং \quad ঘ. ব্লকচেইন মাইনিং }
\item\chTwoMCQBlock{3G-এর সঙ্গে কোন সেবার উন্নয়ন বেশি সম্পর্কিত?
\par ক. মোবাইল ইন্টারনেট ও ভিডিও কল \quad খ. শুধু অ্যানালগ ভয়েস \par গ. শুধু টিভি রিমোট \quad ঘ. কোঅ্যাক্সিয়াল কেবল }
\item\chTwoMCQBlock{4G-এর একটি বৈশিষ্ট্য কোনটি?
\par ক. উচ্চগতির মোবাইল ব্রডব্যান্ড \quad খ. শুধু অ্যানালগ কল \par গ. ইন্টারনেট নেই \quad ঘ. শুধু অবলোহিত যোগাযোগ }
\item\chTwoMCQBlock{5G-এর একটি গুরুত্বপূর্ণ বৈশিষ্ট্য কোনটি?
\par ক. স্বল্প বিলম্ব \quad খ. ডেটা পাঠানো যায় না \par গ. শুধু SMS \quad ঘ. ডিভাইস সংযুক্ত করা যায় না }
\item\chTwoMCQBlock{চলাচলের সময় কল সচল রেখে সেবাদানকারী বেস স্টেশন বদলানোর প্রক্রিয়া কী?
\par ক. Roaming \quad খ. Handoff \par গ. Formatting \quad ঘ. Encoding }
\item\chTwoMCQBlock{নিজস্ব নেটওয়ার্কের বাইরে অন্য সমর্থিত নেটওয়ার্কের সেবা নেওয়াকে কী বলে?
\par ক. Handoff \quad খ. Roaming \par গ. ফোন বন্ধ করা \quad ঘ. আপলোড করা }
\end{enumerate}\end{multicols}
\textbf{উত্তরমালা: ১. খ, ২. খ, ৩. খ, ৪. খ, ৫. ক, ৬. ক, ৭. ক, ৮. ক, ৯. খ, ১০. খ}
\endgroup
